\clearpage
\phantomsection
\addcontentsline{toc}{chapter}{Simulado III — Preparatório AV2 (20\%) e Engenharia Financeira}
\chapter*{Simulado III — Estudo Dirigido Quantitativo (AV2 — 20\%)}
\markboth{SIMULADO III — PREPARATÓRIO AV2}{SIMULADO III — PREPARATÓRIO AV2}

\noindent\tagAV{SIMULADO AV2 (20\%)}\ \tagQuant\ \tagTerritorio
\vspace{6pt}

\begin{tcolorbox}[colback=bgSoft, colframe=etIndigo, boxrule=0.8pt, arc=3pt, title={\bfseries INSTRUÇÕES DE APLICAÇÃO E RESOLUÇÃO}, coltitle=white, colbacktitle=etIndigo]
Este simulado reproduz o padrão quantitativo e analítico da \textbf{AV2 (20\%)} e consolida a engenharia financeira exigida na defesa da \textbf{AV3}. Resolva as \textbf{8 questões objetivas} (múltipla escolha A--E) e as \textbf{2 questões dissertativas} sem consultar o gabarito antes de concluir todas as contas. Ao final, confira cada passo numérico no Gabarito Integralmente Comentado e preencha a Tabela de Diagnóstico.
\end{tcolorbox}

\section*{Parte I — Questões Objetivas de Cálculo Integral (8 Questões)}

\begin{enumerate}[leftmargin=1.6em, itemsep=8pt]

  \item \textbf{(Matriz Locacional Ponderada — Cap. 12)} Uma agroindústria de bioembalagens avalia três municípios ($A$, $B$ e $C$) segundo três critérios ponderados: Logística ($w_1 = 0{,}40$), Infraestrutura Energética/Hídrica ($w_2 = 0{,}35$) e Proximidade a Centros de P\&D ($w_3 = 0{,}25$). As notas atribuídas de $0$ a $10$ foram:
  \begin{center}
  \small
  \begin{tabular}{@{}lccc@{}}
    \toprule
    \textbf{Critério (Peso)} & \textbf{Município A} & \textbf{Município B} & \textbf{Município C} \\
    \midrule
    Logística ($0{,}40$) & 8 & 9 & 7 \\
    Energia e Água ($0{,}35$) & 8 & 6 & 8 \\
    Centros de P\&D ($0{,}25$) & 6 & 8 & 8 \\
    \bottomrule
  \end{tabular}
  \end{center}
  Com base na Matriz Locacional Ponderada, o município vencedor e sua pontuação final são:
  \begin{alternativas}
    \item Município A, com 7,50 pontos.
    \item Município C, com 7,60 pontos.
    \item Município B, com 7,70 pontos.
    \item Município B, com 7,85 pontos.
    \item Empate entre B e C, com 7,65 pontos.
  \end{alternativas}

  \item \textbf{(Capital de Giro e Ponto de Equilíbrio — Cap. 12)} Uma startup de sensores ambientais apresenta os seguintes prazos médios operacionais: $PME = 25$ dias, $PMR = 35$ dias e $PMP = 20$ dias, com OPEX desembolsável anual de \reais{900.000} (ano comercial de $360$ dias). Em sua estrutura de custos, os custos fixos anuais somam \reais{144.000}, o preço unitário de venda é \reais{150,00} e o custo variável unitário é \reais{90,00}. Assinale a alternativa que apresenta corretamente a Necessidade Líquida de Capital de Giro ($NLCG$) e o Ponto de Equilíbrio Contábil em unidades ($PE_{\text{unid}}$):
  \begin{alternativas}
    \item $NLCG = \reais{80.000}$ e $PE_{\text{unid}} = 1.600$ unidades.
    \item $NLCG = \reais{100.000}$ e $PE_{\text{unid}} = 2.400$ unidades.
    \item $NLCG = \reais{100.000}$ e $PE_{\text{unid}} = 1.600$ unidades.
    \item $NLCG = \reais{150.000}$ e $PE_{\text{unid}} = 2.400$ unidades.
    \item $NLCG = \reais{200.000}$ e $PE_{\text{unid}} = 2.880$ unidades.
  \end{alternativas}

  \item \textbf{(Valor Presente Líquido — Cap. 13)} Um projeto de irrigação inteligente exige investimento inicial de \reais{120.000} em $t=0$ e gera fluxos de caixa livres de \reais{55.000} no Ano $1$, \reais{60.500} no Ano $2$ e \reais{66.550} no Ano $3$. Considerando uma $TMA = 10\%$ a.a., o Valor Presente Líquido ($VPL$) do projeto é:
  \begin{alternativas}
    \item $+\reais{15.000,00}$
    \item $+\reais{25.000,00}$
    \item $+\reais{30.000,00}$
    \item $+\reais{45.000,00}$
    \item $+\reais{62.050,00}$
  \end{alternativas}

  \item \textbf{(Payback Descontado — Cap. 13)} Uma cooperativa investe \reais{100.000} em $t=0$ para modernizar sua unidade de beneficiamento, projetando fluxos de caixa livres de \reais{44.000} (Ano $1$), \reais{48.400} (Ano $2$) e \reais{53.240} (Ano $3$). Adotando $TMA = 10\%$ a.a., o Payback Descontado ($PBD$) do projeto é de:
  \begin{alternativas}
    \item $2{,}00$ anos.
    \item $2{,}15$ anos.
    \item $2{,}25$ anos.
    \item $2{,}50$ anos.
    \item $2{,}75$ anos.
  \end{alternativas}

  \item \textbf{(VAUE para Vidas Úteis Diferentes — Cap. 13)} Uma empresa deve escolher entre dois equipamentos mutuamente excludentes e renováveis à $TMA = 10\%$ a.a.: o Equipamento Alfa tem vida útil de $2$ anos e $VPL_{\text{Alfa}} = \reais{21.000,00}$; o Equipamento Beta tem vida útil de $3$ anos e $VPL_{\text{Beta}} = \reais{24.868,52}$. Calcule o Valor Anual Uniforme Equivalente ($VAUE$) de cada alternativa e assinale a decisão correta:
  \begin{alternativas}
    \item $VAUE_{\text{Alfa}} = \reais{10.500,00/\text{ano}}$ e $VAUE_{\text{Beta}} = \reais{8.289,51/\text{ano}}$; escolher Alfa.
    \item $VAUE_{\text{Alfa}} = \reais{12.100,00/\text{ano}}$ e $VAUE_{\text{Beta}} = \reais{10.000,00/\text{ano}}$; escolher Alfa.
    \item $VAUE_{\text{Alfa}} = \reais{12.100,00/\text{ano}}$ e $VAUE_{\text{Beta}} = \reais{12.434,26/\text{ano}}$; escolher Beta.
    \item $VAUE_{\text{Alfa}} = \reais{11.550,00/\text{ano}}$ e $VAUE_{\text{Beta}} = \reais{10.000,00/\text{ano}}$; escolher Alfa.
    \item Como $VPL_{\text{Beta}} > VPL_{\text{Alfa}}$, deve-se escolher Beta independentemente da vida útil.
  \end{alternativas}

  \item \textbf{(WACC — Cap. 14)} Uma empresa de base tecnológica possui valor de mercado do capital próprio $E = \reais{6.000.000}$ e dívida onerosa $D = \reais{4.000.000}$. O custo do capital próprio é $K_e = 15\%$ a.a., o custo bruto da dívida é $K_d = 10\%$ a.a. e a alíquota efetiva de IR/CSLL é $T = 34\%$. O Custo Médio Ponderado de Capital ($WACC$) dessa empresa é:
  \begin{alternativas}
    \item $10{,}64\%$ a.a.
    \item $11{,}64\%$ a.a.
    \item $12{,}00\%$ a.a.
    \item $13{,}00\%$ a.a.
    \item $13{,}40\%$ a.a.
  \end{alternativas}

  \item \textbf{(FCD, Gordon, EV e Equity Value — Cap. 14)} Na avaliação por Fluxo de Caixa Descontado de uma empresa com $WACC = 10\%$ a.a., projetaram-se $FCL_1 = \reais{220.000}$ e $FCL_2 = \reais{242.000}$, além de crescimento perpétuo $g = 2\%$ a.a. após o Ano $2$. Sabe-se que a empresa possui Dívida Bruta de \reais{1.000.000} e Caixa de \reais{300.000}. O \textit{Enterprise Value} ($EV$) e o \textit{Equity Value} ($Equity$) são, respectivamente:
  \begin{alternativas}
    \item $EV = \reais{2.550.000}$ e $Equity = \reais{1.850.000}$.
    \item $EV = \reais{2.950.000}$ e $Equity = \reais{2.250.000}$.
    \item $EV = \reais{2.950.000}$ e $Equity = \reais{1.950.000}$.
    \item $EV = \reais{3.085.500}$ e $Equity = \reais{2.385.500}$.
    \item $EV = \reais{3.485.500}$ e $Equity = \reais{2.785.500}$.
  \end{alternativas}

  \item \textbf{(Método de Venture Capital — Cap. 14)} Uma startup de biotecnologia negocia um aporte de \reais{2.000.000} com um fundo de Venture Capital que exige TIR-alvo de $50\%$ a.a. em um horizonte de saída de $3$ anos. O valor de saída estimado da empresa no Ano $3$ (via múltiplo setorial de $EV/EBITDA$) é de \reais{27.000.000}. Pelo Método de Venture Capital, o \textit{Post-money Valuation}, o \textit{Pre-money Valuation} e o percentual de participação societária exigido pelo fundo são, respectivamente:
  \begin{alternativas}
    \item \reais{8.000.000}; \reais{6.000.000}; $25{,}00\%$.
    \item \reais{9.000.000}; \reais{7.000.000}; $22{,}22\%$.
    \item \reais{6.000.000}; \reais{4.000.000}; $33{,}33\%$.
    \item \reais{10.000.000}; \reais{8.000.000}; $20{,}00\%$.
    \item \reais{8.000.000}; \reais{8.000.000}; $25{,}00\%$.
  \end{alternativas}

\end{enumerate}

\section*{Parte II — Questões Dissertativas / Estudo Dirigido (2 Questões)}

\begin{enumerate}[leftmargin=1.6em, itemsep=8pt]
  \item \textbf{(Integração Locacional, Ponto de Equilíbrio e Análise de Investimentos — Caps. 12 e 13)} Uma cooperativa regional analisa implantar uma planta de processamento de polpa e biomassa de frutas no interior da Bahia. Explique tecnicamente: (a) como o Índice Material de Weber ($IM$) orienta a escolha entre localizar a planta junto à zona rural produtora ou na capital consumidora; (b) por que a análise do Ponto de Equilíbrio Contábil ($PE$) é necessária para gestão operacional, mas insuficiente para aprovar o investimento; e (c) quais dois problemas metodológicos da TIR justificam o uso conjunto do $VPL$ e da $MTIR$. \espacoResolucao[2.6cm]

  \item \textbf{(Valuation de Startups Tecnológicas e Projetos de Impacto — Cap. 14)} Uma agtech universitária ainda pré-lucrativa desenvolve sensores de baixo custo para agricultura familiar no Semiárido. Explique: (a) por que o Fluxo de Caixa Descontado tradicional sofre com o \textit{Dark Side of Valuation} nesse estágio; (b) como o Método de Venture Capital contorna essa limitação na definição de participação societária; e (c) como o conceito de Opções Reais (por exemplo, implantação modular escalável) agrega valor não capturado pelo $VPL$ estático. \espacoResolucao[2.6cm]
\end{enumerate}

\clearpage
\section*{Gabarito Integralmente Comentado — Simulado III}

\begin{boxGabarito}[Resolução Passo a Passo das Questões Objetivas (1 a 8)]
\begin{enumerate}[leftmargin=1.5em, itemsep=5pt]
  \item \textbf{Resposta: (C) Município B, com 7,70 pontos.} \\
  $Nota_A = (0{,}40 \times 8) + (0{,}35 \times 8) + (0{,}25 \times 6) = 3{,}20 + 2{,}80 + 1{,}50 = 7{,}50$. \\
  $Nota_B = (0{,}40 \times 9) + (0{,}35 \times 6) + (0{,}25 \times 8) = 3{,}60 + 2{,}10 + 2{,}00 = \mathbf{7{,}70}$. \\
  $Nota_C = (0{,}40 \times 7) + (0{,}35 \times 8) + (0{,}25 \times 8) = 2{,}80 + 2{,}80 + 2{,}00 = 7{,}60$.

  \item \textbf{Resposta: (B) $NLCG = \reais{100.000}$ e $PE_{\text{unid}} = 2.400$ unidades.} \\
  Ciclo Financeiro: $CF_{ciclo} = 25 + 35 - 20 = 40\text{ dias}$. \\
  Desembolso diário $= 900.000 / 360 = \reais{2.500/\text{dia}}$. \\
  $NLCG = 40 \times 2.500 = \mathbf{\reais{100.000}}$. \\
  Margem de Contribuição unitária: $MC_u = 150 - 90 = \reais{60{,}00}$. \\
  $PE_{\text{unid}} = 144.000 / 60 = \mathbf{2.400\text{ unidades}}$.

  \item \textbf{Resposta: (C) $+\reais{30.000,00}$.} \\
  $VP_1 = 55.000 / 1{,}10 = \reais{50.000{,}00}$. \\
  $VP_2 = 60.500 / 1{,}10^2 = 60.500 / 1{,}21 = \reais{50.000{,}00}$. \\
  $VP_3 = 66.550 / 1{,}10^3 = 66.550 / 1{,}331 = \reais{50.000{,}00}$. \\
  $VPL = -120.000 + 50.000 + 50.000 + 50.000 = \mathbf{+\reais{30.000{,}00}}$.

  \item \textbf{Resposta: (D) $2{,}50$ anos.} \\
  $VP_1 = 44.000 / 1{,}10 = \reais{40.000}$ $\Rightarrow$ Saldo descontado $= -100.000 + 40.000 = -\reais{60.000}$. \\
  $VP_2 = 48.400 / 1{,}21 = \reais{40.000}$ $\Rightarrow$ Saldo descontado $= -60.000 + 40.000 = -\reais{20.000}$. \\
  $VP_3 = 53.240 / 1{,}331 = \reais{40.000}$. \\
  $PBD = 2 + \frac{20.000}{40.000} = \mathbf{2{,}50\text{ anos}}$.

  \item \textbf{Resposta: (B) $VAUE_{\text{Alfa}} = \reais{12.100,00/\text{ano}}$ e $VAUE_{\text{Beta}} = \reais{10.000,00/\text{ano}}$; escolher Alfa.} \\
  $FRC(10\%, 2) = \frac{0{,}10 \times 1{,}21}{1{,}21 - 1} = \frac{0{,}121}{0{,}21} = 0{,}57619048 \Rightarrow VAUE_{\text{Alfa}} = 21.000 \times 0{,}57619048 = \mathbf{\reais{12.100{,}00/\text{ano}}}$. \\
  $FRC(10\%, 3) = \frac{0{,}10 \times 1{,}331}{1{,}331 - 1} = \frac{0{,}1331}{0{,}331} = 0{,}40211480 \Rightarrow VAUE_{\text{Beta}} = 24.868{,}52 \times 0{,}40211480 = \mathbf{\reais{10.000{,}00/\text{ano}}}$. \\
  Como $VAUE_{\text{Alfa}} > VAUE_{\text{Beta}}$, escolhe-se \textbf{Alfa}.

  \item \textbf{Resposta: (B) $11{,}64\%$ a.a.} \\
  $D + E = 4.000.000 + 6.000.000 = \reais{10.000.000}$ ($w_E = 0{,}60$, $w_D = 0{,}40$). \\
  $WACC = (0{,}60 \times 15\%) + (0{,}40 \times 10\% \times (1 - 0{,}34)) = 9{,}00\% + 2{,}64\% = \mathbf{11{,}64\%\text{ a.a.}}$.

  \item \textbf{Resposta: (B) $EV = \reais{2.950.000}$ e $Equity = \reais{2.250.000}$.} \\
  $VP(FCL_1) = 220.000 / 1{,}10 = \reais{200.000}$. \\
  $VP(FCL_2) = 242.000 / 1{,}21 = \reais{200.000}$. \\
  $VT_2 = \frac{242.000 \times 1{,}02}{0{,}10 - 0{,}02} = \frac{246.840}{0{,}08} = \reais{3.085.500}$. \\
  $VP(VT_2) = 3.085.500 / 1{,}21 = \mathbf{\reais{2.550.000}}$. \\
  $EV = 200.000 + 200.000 + 2.550.000 = \mathbf{\reais{2.950.000}}$. \\
  $Equity = 2.950.000 - 1.000.000 + 300.000 = \mathbf{\reais{2.250.000}}$.

  \item \textbf{Resposta: (A) \reais{8.000.000}; \reais{6.000.000}; $25{,}00\%$.} \\
  $\text{Post-money} = \frac{27.000.000}{(1{,}50)^3} = \frac{27.000.000}{3{,}375} = \mathbf{\reais{8.000.000}}$. \\
  $\text{Pre-money} = 8.000.000 - 2.000.000 = \mathbf{\reais{6.000.000}}$. \\
  $\%\text{Participação} = \frac{2.000.000}{8.000.000} \times 100\% = \mathbf{25{,}00\%}$.
\end{enumerate}
\end{boxGabarito}

\begin{boxGabarito}[Espelho de Correção — Questões Dissertativas (1 e 2)]
\textbf{Questão Dissertativa 1 (10 pontos):}
\begin{itemize}[leftmargin=1.4em, itemsep=2pt]
  \item \textbf{(a) Índice Material de Weber (3,5 pts):} Explicar que $IM = \text{Peso MP}/\text{Peso PA} > 1$ indica grande perda de massa (casca, água) no beneficiamento de frutas, tornando o frete da matéria-prima bruta mais caro e orientando a planta para a zona produtora rural.
  \item \textbf{(b) Limite do Ponto de Equilíbrio (3,0 pts):} Apontar que o $PE$ indica apenas o volume em que o lucro contábil operacional zera ($LAJIR = 0$), mas ignora o custo de oportunidade do capital ($TMA$), o valor do dinheiro no tempo e o investimento inicial ($CAPEX + NLCG$).
  \item \textbf{(c) Problemas da TIR vs. VPL e MTIR (3,5 pts):} Citar (i) múltiplas TIRs ou ausência de raiz em fluxos não convencionais (inversões de sinal) e (ii) hipótese irrealista de reinvestimento dos fluxos intermediários à própria TIR; explicar que o $VPL$ mede riqueza absoluta à TMA e a $MTIR$ corrige o reinvestimento usando $k_r = TMA$.
\end{itemize}
\vspace{4pt}
\textbf{Questão Dissertativa 2 (10 pontos):}
\begin{itemize}[leftmargin=1.4em, itemsep=2pt]
  \item \textbf{(a) Dark Side of Valuation (3,5 pts):} Explicar a ausência de histórico financeiro, fluxos de caixa negativos iniciais, alta incerteza de adoção tecnológica e dependência excessiva do Valor Terminal no FCD.
  \item \textbf{(b) Método de Venture Capital (3,5 pts):} Descrever o cálculo retroativo a partir do valor de saída futuro (múltiplo sobre métrica projetada), descontado pela TIR-alvo do fundo para definir \textit{Post-money}, \textit{Pre-money} e percentual societário.
  \item \textbf{(c) Opções Reais (3,0 pts):} Explicar que o $VPL$ estático trata a decisão como irreversível e rígida, enquanto Opções Reais precificam a flexibilidade de adiar, expandir em módulos após validação ou abandonar limitando perdas ($VPL_{\text{expandido}} = VPL_{\text{estático}} + VOR$).
\end{itemize}
\end{boxGabarito}

\section*{Tabela de Diagnóstico para a AV2 (Autoavaliação)}

\renewcommand{\arraystretch}{1.25}
\small
\begin{center}
\begin{tabular}{@{}p{4.6cm} c p{5.5cm} p{2.8cm}@{}}
  \toprule
  \rowcolor{etIndigo}
  \textcolor{white}{\bfseries Eixo Avaliado} & \textcolor{white}{\bfseries Questões} & \textcolor{white}{\bfseries Seu Desempenho (Acertos / Notas)} & \textcolor{white}{\bfseries Revisão Indicada} \\
  \midrule
  Localização, NLCG e Ponto de Equilíbrio & Obj. 1, 2 + Diss. 1 & \rule{2.5cm}{0.4pt} / 2 obj. + \rule{1.2cm}{0.4pt} pts & Cap. 12 (Etapas 3 e 4) \\
  \rowcolor{bgGray}
  VPL, Payback Descontado e VAUE & Obj. 3, 4, 5 + Diss. 1 & \rule{2.5cm}{0.4pt} / 3 obj. + \rule{1.2cm}{0.4pt} pts & Cap. 13 (Etapas 3 e 4) \\
  WACC, Gordon, EV/Equity e Venture Capital & Obj. 6, 7, 8 + Diss. 2 & \rule{2.5cm}{0.4pt} / 3 obj. + \rule{1.2cm}{0.4pt} pts & Cap. 14 (Etapas 3 e 4) \\
  \bottomrule
\end{tabular}
\end{center}
\normalsize
